\section{Conclusion}
\label{sec:conclusion}

We presented \method{}, a training-free, fixed-budget cache-management framework that turns a causal streaming visual geometry transformer into a continual perception module for robotic manipulation. \method{} retains the cached tokens that later viewpoints depend on, as measured by cross-frame attention saliency, stabilized spatially and temporally, and balanced against feature diversity under budgeted per-layer allocation, so that memory and per-step compute stay constant regardless of stream length. On single-camera DROID manipulation streams, \method{} matches the reconstruction quality of the strongest memory-bounded baselines, roughly doubles their throughput, and holds peak VRAM near 9.5\,GB where the full-cache backbone fails at 300 frames. Because geometry quality is nearly flat across a $16\times$ budget range, the token budget serves as a deployment knob that changes cost, not accuracy. Limitations remain, as causal inference cannot retrospectively correct early errors and the evaluation relies on model-generated reference labels from a single camera. Future work includes multi-camera fusion, lightweight anchor re-alignment, and closing the loop with manipulation policies and 3D world models.
